% ============================================================
\chapter{Research Publication}
\label{chap:pub}
% ============================================================

\section[\projecttitle]{\raggedright\projecttitle}

A research paper drawn from this work, ``\projecttitle'', is written in the IEEE
conference format by \candone, \candtwo\ and \candthree, with \guidename\ and
\mentorname\ of the \department, \institute. Status: \publicationstatus.

The paper presents PipedPeer for a systems readership and is organised as
follows.

\begin{itemize}[leftmargin=*]
  \item \textbf{Abstract and Introduction} state the three obstacles to using
  nearby idle machines --- the environment, the code and the network --- and the
  paper's contributions.
  \item \textbf{Related Work} places PipedPeer against cluster frameworks for
  Python, volunteer computing, environment reproduction, NAT traversal, and
  communication-efficient data-parallel training.
  \item \textbf{Design} sets out the design goals and the architecture of
  Chapter~\ref{chap:arch}: one daemon per machine, coordination by the submitting
  machine, and an introducer in every daemon.
  \item \textbf{Implementation} summarises the mechanisms of
  Chapter~\ref{chap:design}: content-addressed environments and path-level
  transfer, transparent interception with measured work sharing and the cost
  model, training synchronised through the daemon, and joining across networks.
  \item \textbf{Evaluation} states its questions and answers them with the
  measurements of Chapter~\ref{chap:results}, in the order of setup,
  microbenchmarks, end-to-end runs and training.
  \item \textbf{Limitations and Conclusion} close the paper.
\end{itemize}
